\chapter{Literature Review}\label{chap:literature}

\section{Introduction}

\ProjectName{} sits at the intersection of several fields: emergency medical services, decision support and
triage, software architecture, real-time communication, geospatial computing, and privacy and security
engineering. This chapter reviews the work in each field that shaped the design, explains the choices that follow
from it, and closes by identifying the gap that the platform addresses. The review is selective: it concentrates
on the ideas that the system actually uses.

\section{Emergency Dispatch and Emergency Medical Services}

An emergency response is commonly described as a sequence of stages: detection and reporting, call-taking,
triage, dispatch, travel, on-scene care, transport and hospital handover. Computer-aided dispatch (CAD) systems
support the middle stages by recording incidents, tracking resources and recommending assignments. Their
quality depends on the information that reaches them: the location, the nature of the emergency and its urgency.

A large body of operations-research literature studies how the \emph{resources} of this chain should be
arranged. Location models decide where to base ambulances so that a target share of the population can be
reached within a time threshold; relocation models decide how to move vehicles as demand and availability change.
Gendreau, Laporte and Semet formulated a double-coverage ambulance location model and solved it with tabu search
\cite{gendreau1997tabu}; Brotcorne, Laporte and Semet reviewed location and relocation models and the dynamic
dispatching problems built on them \cite{brotcorne2003ambulance}. A more recent review by Aringhieri
\emph{et al.}\ widens the scope from ambulance location to the wider emergency-care system and the new
challenges it faces \cite{aringhieri2017ems}.

Two observations from this literature guide \ProjectName{}. First, optimisation models take the demand and the
state of the fleet as given; they presuppose accurate, timely data about where incidents are and which units are
free. A platform that captures location automatically and tracks unit status and position in real time produces
exactly that data. Second, the human dispatcher remains central: the models inform, they do not replace. \ProjectName{}
follows the same philosophy by ranking units for the dispatcher while leaving the assignment to a person.

\section{Triage and Decision Support}

Triage orders patients, or incidents, by urgency. Emergency departments often use structured five-level
instruments such as the Emergency Severity Index \cite{gilboy2012esi}, which combines the patient's acuity
with the resources they are expected to need. Triage at the \emph{report} stage is harder, because the
information is a short free-text message, a category and a location rather than a clinical assessment.

Language models offer a practical way to classify such messages. Large pre-trained models can perform new
classification tasks from instructions and a few examples without task-specific training
\cite{brown2020gpt3}, which suits a domain where labelled emergency data is scarce and sensitive.
\ProjectName{} therefore uses a language model, called through a hosted gateway, to suggest a priority and a
possible duplicate for each report. The design treats the model as an \emph{advisor with a safety net}: the output is a
suggestion shown to a dispatcher; the call is time-limited and deterministic in configuration; and a failure
leaves the report untouched (\cref{sec:ai-components}).

For structured, tabular signals such as vital signs, conventional supervised learning is more appropriate.
Gradient-boosted trees \cite{chen2016xgboost} and the scikit-learn ecosystem \cite{pedregosa2011sklearn} are
well-established tools for such data, and they are the formats the platform's AI model service is built to host
and serve with a strict input contract.

\Needspace{16\baselineskip}
\section{Software Architecture for Emergency Platforms}

\subsection{Modularity}
Parnas's classic argument is that a system should be decomposed into modules according to design decisions
likely to change, each module hiding one such decision \cite{parnas1972decomposing}. Architectural qualities
such as modifiability, availability and security are achieved by deliberate structural choices
\cite{bass2021software}, and enterprise-application patterns such as layered services, repositories and
domain models give those choices a shared vocabulary \cite{fowler2002peaa,evans2003ddd}.

\subsection{Monolith or microservices?}
Microservices decompose an application into independently deployable services that communicate over the network
\cite{lewis2014microservices,newman2021microservices,richardson2018microservices}. They offer independent
scaling and release, but they also impose a cost in distributed-systems complexity: network failure, data
consistency across services, operational overhead and the need for substantial infrastructure. A single
deployable application organised internally into well-separated modules --- a \emph{modular monolith} --- keeps
most of the benefits of modularity while avoiding that cost, and its modules can be extracted later if a real
need arises. Newman himself discusses the merits of starting a new system as a monolith and splitting it later
\cite{newman2021microservices}.

\ProjectName{} follows this reasoning. For a team of this size, a modular Spring Boot backend, organised by feature
module with consistent controller--service--repository layering, gives strong boundaries, simple transactions and a
simple deployment. The one component that genuinely differs in nature --- the AI model service, written in Python and
evolving with machine-learning tooling --- is deployed separately and reached over an authenticated interface
(\cref{chap:design}).

\subsection{Architecture of the clients}
The two clients follow current practice for their platforms: a declarative, state-driven interface in Flutter
\cite{flutter,riverpod} and the React-based App Router of Next.js with server-state caching \cite{nextjs,tanstackquery}.

\section{Real-Time and Event-Driven Communication}

Emergency software must show changes as they happen. Polling wastes bandwidth and adds delay; a persistent
bidirectional channel is better. The WebSocket protocol provides one over a single TCP connection
\cite{rfc6455}, and STOMP adds a simple publish--subscribe messaging layer on top of it \cite{stomp12}.
\ProjectName{} uses STOMP over WebSocket for live updates and Firebase Cloud Messaging for push notifications to
people who are away from the app \cite{firebase}.

Integration patterns for messaging \cite{hohpe2003eip} and for data-intensive systems \cite{kleppmann2017ddia}
motivate attention to a subtle hazard: if a message is published from inside a database transaction, a consumer can react
before the data is committed and read stale state; if the transaction later rolls back, the message announces
something that never happened. The design therefore publishes live updates and notifications only \emph{after the
transaction commits} (\cref{sec:after-commit}).

\section{Geospatial Services}

Locating incidents and units requires geographic data and distance computation. The great-circle distance between
two coordinates is conveniently computed with the haversine formula \cite{sinnott1984haversine}, which is accurate
enough to rank nearby units. For real travel times, a routing engine is needed. OpenStreetMap is a
community-maintained global map database \cite{haklay2008osm}, and the Open Source Routing Machine computes fast
routes on OpenStreetMap data \cite{luxen2011osrm}. Reverse geocoding turns coordinates into a readable address
\cite{nominatim}, and vector map rendering is provided by MapLibre \cite{maplibre}. \ProjectName{} combines
these: the console ranks units by distance and estimated travel time, and both clients draw routes and maps.

\section{Security, Privacy and Trust}

\subsection{Design principles}
Saltzer and Schroeder's design principles for protection remain the baseline: least privilege, fail-safe
defaults, complete mediation and psychological acceptability \cite{saltzer1975protection}. The OWASP Top 10
catalogues the most common web-application weaknesses, among them broken access control, cryptographic failures
and injection \cite{owasp2021top10}. In \ProjectName{} every request is authenticated, authorised by role and by
ownership inside the service layer, and audited where it touches sensitive data.

\subsection{Authentication and secrets}
Signed tokens such as JSON Web Tokens \cite{rfc7519} let a server verify a caller without a session store; the
platform delegates credential handling to a managed identity provider \cite{firebase}, so passwords never reach its
own code. One-time codes and tokens are stored only as keyed hashes (HMAC) \cite{rfc2104}, so a leaked database does
not reveal a usable code.

\subsection{Privacy of health data}
A medical profile is among the most sensitive data a system can hold. Egyptian Law No.\ 151 of 2020 on personal
data protection \cite{egypt2020pdpl} establishes principles of lawful processing, purpose limitation and data-subject
rights. The platform applies them as engineering rules: sharing is based on explicit, withdrawable consent; access
is limited to those who are treating the patient and only while the incident is open; free-text medical fields are
encrypted at rest; and every access is recorded.

\subsection{Abuse control without denial}
Rate limiting is commonly implemented with the HTTP status \texttt{429 Too Many Requests} and a
\texttt{Retry-After} hint \cite{rfc6585}. For an emergency line, outright refusal is unacceptable, so the caller-trust policy
only \emph{slows} repeat abusers and always tells them which emergency numbers to call (\cref{sec:caller-trust}).

\section{Reliability Patterns}

Production systems fail in partial and unexpected ways. Nygard catalogues stability patterns --- time-outs, circuit
breakers, bulkheads and graceful degradation --- that keep a failure of one dependency from taking down the
whole system \cite{nygard2018release}. Idempotent operations make retries safe, which matters on unreliable mobile
networks \cite{kleppmann2017ddia}. \ProjectName{} applies these ideas directly: each external service is wrapped so its
failure disables only its own feature, report submission is idempotent through a client-generated request id, and
the mobile app keeps an offline outbox.

\section{Related Work}

Rather than naming individual products, whose capabilities change and are not always publicly documented, this
section compares the \emph{categories} of system that a team would encounter when looking for an existing solution
(\cref{tab:related}). The comparison is qualitative and describes typical characteristics; individual systems vary.

\begin{table}[!htbp]
\caption{Typical characteristics of categories of emergency-related systems compared with \ProjectName{}.}\label{tab:related}
\centering\tabfont\setlength{\tabcolsep}{4.5pt}\rowband
\begin{tabular}{@{}>{\raggedright\arraybackslash}p{4.9cm}*{4}{>{\centering\arraybackslash}p{2.2cm}}@{}}
\rowcolor{najdanavy}\textcolor{white}{\sffamily\bfseries\small Capability} &
  \textcolor{white}{\sffamily\bfseries\scriptsize Call-taking and CAD} &
  \textcolor{white}{\sffamily\bfseries\scriptsize Citizen SOS apps} &
  \textcolor{white}{\sffamily\bfseries\scriptsize Research prototypes} &
  \textcolor{white}{\sffamily\bfseries\scriptsize \ProjectName} \\
Citizen self-service reporting with live location & partial & \yes & varies & \yes \\
Evidence attachments (photo, video, audio) & rare & partial & varies & \yes \\
Automated priority triage and duplicate hints & rare & \no & varies & \yes \\
Live dispatcher queue and unit assignment & \yes & \no & partial & \yes \\
Field-crew mobile mission workflow & partial & \no & partial & \yes \\
Hospital pre-notification with live vitals & rare & \no & rare & \yes \\
Consent-based, audited medical-profile sharing & rare & partial & rare & \yes \\
Graduated abuse control that never refuses a caller & rare & rare & \no & \yes \\
Arabic and English with right-to-left layout & varies & varies & varies & \yes \\
\end{tabular}
\end{table}

\section{Research Gap and Position of This Work}

The literature offers strong models for allocating emergency resources, good tools for classification, and mature
engineering patterns for reliable, secure systems. What it rarely offers is an \emph{integrated, working platform}
that closes the loop between the citizen who reports, the dispatcher who decides, the crew who responds and the
hospital that receives --- and that treats privacy, consent and abuse as part of the design rather than as afterthoughts.
\ProjectName{} fills that gap for the Egyptian setting: it combines automatic location and AI-assisted triage with a
live dispatch workflow, a consent-based medical profile, hospital pre-notification and a fail-open abuse policy, in
Arabic and English, using an architecture proportionate to a small team.

\section{Chapter Summary}

The review established the principles behind the design: human-in-the-loop decision support, a modular
architecture rather than premature distribution, real-time communication with after-commit event publication,
standard geospatial tooling, least-privilege security with strong privacy for health data, and reliability
patterns for partial failure. The next chapter turns these principles into a concrete design.
